% ======================================================================
%  KAVRAM: AI ≠ Traditional Software  (Tier 1)
%  Kitabın açılış kavramı — olasılıksal düşünmenin zemini
% ======================================================================

\concept{Temel · Foundations}{AI Geleneksel\\Yazılım Değildir}{Kod kuralları \emph{icra} eder; model bir sonucu \emph{tahmin} eder}

% ----------------------------------------------------------------------
\stage{1}{Problem}
% ----------------------------------------------------------------------
Bir mühendis ilk kez bir LLM özelliği yazdığında, onu eski refleksleriyle
karşılar: aynı girdiyi iki kez verir, iki \emph{farklı} cevap alır. Dün geçen
testi bugün kırılır. ``Bug yok ki --- kod aynı'' der ve şaşırır.

\keyline{Sorun kodun bozulması değil. Sorun, deterministik bir araca davranır gibi olasılıksal bir araçla çalışıyor olman.}

Geleneksel yazılımda bir fonksiyon bir \emph{sözleşmedir}: \texttt{topla(2,3)}
her zaman \texttt{5} döner. LLM ise sözleşme değil, \emph{olasılık} sunar: ``bu
girdiye en olası devam şudur.'' Bu farkı kavramadan kurulan her sistem,
güvenilirliği yanlış yerde arar --- daha çok test, daha büyük model, daha katı
prompt --- ve hiçbiri kök nedeni çözmez.

% ----------------------------------------------------------------------
\stage{2}{Neden Var}
% ----------------------------------------------------------------------
Fark, iki yazılımın \textbf{nasıl üretildiğinden} doğar. Geleneksel kodu bir
insan \emph{yazar}: davranış, satır satır açıkça programlanır. Bir modeli ise
kimse yazmaz; model milyonlarca örnekten \emph{istatistiksel olarak eğitilir}.
Davranışı kurallarda değil, ağırlıklarda --- bir olasılık dağılımında --- saklıdır.

Bu yüzden:
\begin{itemize}
  \item Davranışı ``düzeltmek'' bir satırı değiştirmek değil, girdiyi/bağlamı/kısıtı değiştirmektir.
  \item ``Doğruluk'' ikili (doğru/yanlış) değil, \emph{olasılıksaldır} (genelde doğru).
  \item Test etmek, tek bir çıktıyı değil, çıktı \emph{dağılımını} doğrulamaktır.
\end{itemize}

% ----------------------------------------------------------------------
\stage{3}{Tarihçe}
% ----------------------------------------------------------------------
\begin{timeline}
\yr{1960–80} kural tabanlı / uzman sistemler (elle \texttt{if-then}, deterministik) \tlsep \yr{1990–2010} makine öğrenmesi: davranış veriden öğrenilir \tlsep \yr{2012–17} derin öğrenme: özellikler de öğrenilir \tlsep \yr{2018–26} LLM'ler: ``programlama'' → ``yönlendirme'' (prompt + context)
\end{timeline}

% ----------------------------------------------------------------------
\stage{4}{Tanım}
% ----------------------------------------------------------------------
\begin{problembox}
\textbf{Deterministik sistem:} Aynı girdi, her zaman aynı çıktıyı verir;
davranış açıkça programlanmıştır (\emph{execution} --- icra).

\smallskip
\textbf{Olasılıksal sistem:} Aynı girdi, bir olasılık dağılımından
\emph{örneklenmiş} çıktılar verir; davranış eğitimle öğrenilmiştir
(\emph{prediction} --- tahmin).
\end{problembox}

LLM ikincisidir: her adımda ``bir sonraki en olası token''ı bir dağılımdan seçer.
\term{temperature} bu seçimin ne kadar çeşitli olacağını ayarlar --- ama sıfırda
bile model bir \emph{sözleşme} değil, bir \emph{tahmin} üretir.

% ----------------------------------------------------------------------
\stage{5}{Zihinsel Model}
% ----------------------------------------------------------------------
\begin{mentalbox}
Geleneksel yazılım bir \textbf{otomat} gibidir: B4'e basarsın, her seferinde
\emph{aynı} cipsi verir. Garanti, mekaniktir.

\smallskip
AI ise \textbf{çok okumuş, çok hızlı bir stajyer} gibidir: aynı soruyu iki kez
sorarsan, ikisi de doğru ama \emph{farklı kelimelerle} iki cevap alırsın ---
ve bazen büyük bir özgüvenle yanlış cevap verir. Onu \emph{programlamazsın};
ona iyi bir brief verir, işini kontrol eder ve sınırlarını çizersin.
\end{mentalbox}

Bu model neden önemli: bir stajyerden ``her seferinde birebir aynı paragrafı''
beklemezsin. Beklersen sorun stajyerde değil, beklentindedir.

% ----------------------------------------------------------------------
\stage{6}{Kavramsal Diyagram}
% ----------------------------------------------------------------------
Aynı girdi, iki farklı dünya. Solda tek ve kesin çıktı; sağda bir dağılımdan
örneklenen çıktı.

\vspace{4pt}
\begin{center}
\begin{tikzpicture}[
   font=\sffamily\scriptsize,
   io/.style={draw=softgray, line width=0.7pt, fill=white, rounded corners=8pt,
              minimum height=8mm, minimum width=20mm, align=center, inner sep=3pt},
   proc/.style={draw=accent, line width=0.9pt, fill=wash, rounded corners=2pt,
                minimum height=11mm, minimum width=26mm, align=center},
   outbox/.style={draw=good, line width=0.9pt, fill=good!10, rounded corners=2pt,
               minimum height=8mm, minimum width=20mm, align=center},
   flow/.style={-{Stealth[length=2mm]}, accent, line width=0.8pt}]

  % --- SOL: Geleneksel ---
  \node[io] (in1) {Girdi};
  \node[proc, below=5mm of in1] (rule) {Açık kurallar\\\texttt{if / else}};
  \node[outbox, below=5mm of rule] (o1) {\textbf{Çıktı}\\tek · kesin};
  \draw[flow] (in1) -- (rule);
  \draw[flow] (rule) -- (o1);
  \node[above=2mm of in1, text=ink, font=\sffamily\bfseries]
        {GELENEKSEL · deterministik};

  % ayraç
  \draw[hairline, line width=0.8pt, dotted] ($(in1.north)!0.5!(in1.north)+(33mm,4mm)$)
        -- ++(0,-42mm);

  % --- SAĞ: AI ---
  \node[io, right=46mm of in1] (in2) {Girdi};
  \node[proc, below=5mm of in2, draw=accentdk] (dist) {Olasılık\\dağılımı};
  \draw[flow] (in2) -- (dist);
  \node[above=2mm of in2, text=ink, font=\sffamily\bfseries]
        {AI · olasılıksal};

  % dağılım çubukları
  \begin{scope}[shift={($(dist.south)+(-12mm,-6mm)$)}]
    \def\bw{3.2mm}
    \fill[accent]   (0,0)        rectangle ++(\bw,9mm);
    \fill[accent!55](4.6mm,0)    rectangle ++(\bw,6mm);
    \fill[accent!35](9.2mm,0)    rectangle ++(\bw,4mm);
    \fill[accent!22](13.8mm,0)   rectangle ++(\bw,2.5mm);
    \fill[accent!15](18.4mm,0)   rectangle ++(\bw,1.5mm);
    \draw[hairline] (-1mm,0) -- (23mm,0);
    \node[font=\sffamily\tiny, text=softgray, below=0pt] at (11mm,-1mm)
          {olası çıktılar};
  \end{scope}

  \node[outbox, below=24mm of dist] (o2) {\textbf{Çıktı}\\örneklenmiş};
  \draw[flow] ($(dist.south)+(-10.4mm,-6mm)$) .. controls +(0,-5mm) and +(0,5mm) ..
        node[right=1mm, text=accentdk, font=\sffamily\tiny]{örnekle} (o2.north);
\end{tikzpicture}
\end{center}

\keyline{Solda ``test = tek doğru cevabı kontrol et''. Sağda ``test = dağılımın kabul edilebilir olduğunu kontrol et''.}

% ----------------------------------------------------------------------
\stage{7}{Karar Çerçevesi}
% ----------------------------------------------------------------------
``Bu işi AI mı yapsın, geleneksel kod mu?'' Çoğu üretim hatası bu soruyu
atlamaktan doğar. AI güçlü diye her yere konmaz.

\vspace{2pt}
\begin{center}
\begin{tikzpicture}[
   font=\sffamily\scriptsize, node distance=5mm and 8mm,
   q/.style={draw=accent, line width=0.8pt, fill=wash, diamond, aspect=2.4,
             align=center, inner sep=1pt, minimum width=32mm},
   leaf/.style={draw=accentdk, line width=0.8pt, rounded corners=2pt,
                fill=accent!12, align=center, inner sep=4pt, minimum width=24mm},
   codeleaf/.style={draw=softgray, line width=0.8pt, rounded corners=2pt,
                fill=white, align=center, inner sep=4pt, minimum width=24mm},
   yes/.style={good, font=\sffamily\bfseries\scriptsize},
   no/.style={warn, font=\sffamily\bfseries\scriptsize},
   ed/.style={-{Stealth[length=2mm]}, softgray, line width=0.7pt}]

  % dikey omurga: sorular aşağı iner, yapraklar yana dallanır
  \node[q] (q1) {Kuralı açıkça\\yazabiliyor musun?};
  \node[codeleaf, right=13mm of q1] (code1) {\textbf{Geleneksel kod}\\ucuz · kesin};
  \node[q, below=11mm of q1] (q2) {Girdi muğlak / doğal\\dil / algı mı?};
  \node[codeleaf, left=13mm of q2] (code2) {\textbf{Kod}\\(belki ML sınıflandırıcı)};
  \node[q, below=11mm of q2] (q3) {Çıktıdaki değişkenlik\\kabul edilebilir mi?};
  \node[leaf, below left=12mm and 4mm of q3, minimum width=22mm] (ai) {\textbf{AI}\\+ guardrail + eval};
  \node[codeleaf, below right=12mm and 4mm of q3, minimum width=22mm] (code3) {\textbf{Kısıtlı AI}\\structured output};

  \draw[ed] (q1) -- node[yes,above,inner sep=1pt]{Evet} (code1);
  \draw[ed] (q1) -- node[no,right]{Hayır} (q2);
  \draw[ed] (q2) -- node[no,above,inner sep=1pt]{Hayır} (code2);
  \draw[ed] (q2) -- node[yes,right]{Evet} (q3);
  \draw[ed] (q3) -- node[yes,left]{Evet} (ai);
  \draw[ed] (q3) -- node[no,right]{Hayır} (code3);
\end{tikzpicture}
\end{center}

\keyline{Kural: Açıkça yazılabilen hiçbir şeyi modele yaptırma. AI'ı yalnızca muğlaklığın gerçek olduğu yerde kullan.}

% ----------------------------------------------------------------------
\stage{8}{Üretim Senaryosu}
% ----------------------------------------------------------------------
\begin{prodbox}[Üretim Senaryosu · Kırılgan Test]
\textbf{Bağlam.} Bir ekip, e-postaları özetleyen bir özellik yazar. CI'a klasik
bir test koyar: \texttt{assert ozet == "Bu hafta 3 toplantı var."}

\smallskip
\textbf{Sorun.} Test rastgele kırılır. Model bazen ``Bu hafta üç toplantınız
var.'' der --- anlamca aynı, dizgece farklı. Ekip ``model güvenilmez'' diye
şikâyet eder.

\smallskip
\textbf{Çözüm.} Birebir string yerine \emph{özellik} (property) testi: özet 200
karakterden kısa mı? ``3'' veya ``üç'' geçiyor mu? Bir \term{judge model}
``orijinali doğru özetliyor mu?'' diye puanlıyor mu? Artık test, tek bir çıktıyı
değil \textbf{kabul edilebilir dağılımı} doğrular --- ve kırılgan olmaktan çıkar.
\end{prodbox}

% ----------------------------------------------------------------------
\stage{9}{Anti-pattern}
% ----------------------------------------------------------------------
\begin{antibox}[Anti-pattern · ``AI-First Düşünme'']
Her problemi önce AI ile çözmeye çalışmak. Bir e-posta adresini doğrulamak için
LLM çağırmak (regex yeterken), sabit bir hesaplama için modele güvenmek, ya da
model çıktısını sorgusuz bir \emph{sözleşme} gibi sisteme akıtmak.

\smallskip
\textbf{Belirti:} Basit, deterministik işler pahalı, yavaş ve ara sıra yanlış.\\
\textbf{Kök neden:} Olasılıksal aracı, deterministik işe koşmak.\\
\textbf{Düzeltme:} Önce ``bu kural yazılabilir mi?'' diye sor. Yazılabiliyorsa yaz. AI'ı yalnızca muğlaklık için sakla.
\end{antibox}

% ----------------------------------------------------------------------
\stage{10}{Alternatifler \& Ödünleşimler}
% ----------------------------------------------------------------------
\vspace{2pt}
\begin{tabularx}{\linewidth}{@{}>{\raggedright\arraybackslash}p{30mm} L L@{}}
\arrayrulecolor{hairline}\toprule
{\sffamily\bfseries\small Yaklaşım} &
{\sffamily\bfseries\small Güçlü olduğu yer} &
{\sffamily\bfseries\small Ödünleşim} \\
\midrule
\textbf{Kural tabanlı kod} & Net, sabit, doğrulanabilir mantık &
Muğlaklık ve doğal dilde kırılır \\
\addlinespace[3pt]
\textbf{ML sınıflandırıcı} & Sabit kategoriler, çok örnek var &
Eğitim verisi + bakım; esnek değil \\
\addlinespace[3pt]
\textbf{LLM} & Açık uçlu dil, muğlaklık, az örnek &
Olasılıksal; eval ve guardrail şart \\
\bottomrule
\end{tabularx}

\vspace{4pt}
Pratikte en sağlam sistemler \textbf{hibrittir}: deterministik iskelet (kod) +
olasılıksal çekirdek (model), arada doğrulama katmanıyla.

% ----------------------------------------------------------------------
\stage{11}{İlişkili Kavramlar}
% ----------------------------------------------------------------------
\begin{itemize}
  \item \term{Temperature} --- örneklemedeki çeşitliliği ayarlar; 0 bile determinizm garanti etmez.
  \item \term{Hallucination} --- olasılıksal doğanın doğrudan sonucu: dayanaksız ama emin çıktı.
  \item \term{Eval} --- tek çıktı yerine dağılımı ölçmenin yolu.
  \item \term{Guardrails} --- olasılıksal çıktıyı deterministik sınırlara hapsetme katmanı.
  \item \term{Structured Output} --- modeli sözleşmeye yaklaştıran kısıtlama.
\end{itemize}

% ----------------------------------------------------------------------
\stage{12}{Özet}
% ----------------------------------------------------------------------
\begin{itemize}
  \item Geleneksel yazılım \emph{icra eder} (deterministik); AI \emph{tahmin eder} (olasılıksal).
  \item Fark araçtan değil, üretim biçiminden gelir: kod yazılır, model eğitilir.
  \item ``Doğruluk'' ikili değil olasılıksaldır; test, çıktıyı değil \textbf{dağılımı} doğrular.
  \item Açıkça yazılabilen işi modele yaptırma --- AI'ı muğlaklık için sakla.
  \item En sağlam sistemler hibrittir: deterministik iskelet + olasılıksal çekirdek + doğrulama.
\end{itemize}

% ----------------------------------------------------------------------
\stage{13}{Mühendislik İlkesi}
% ----------------------------------------------------------------------
\begin{principlebox}
\textbf{İlke \#1 — AI'ı programlamazsın; ona brief verir, çıktısını kısıtlar ve dağılımını ölçersin.}\\[4pt]
\normalsize Determinizm beklediğin yerde olasılıksal bir araç kullanma. Model
bir sözleşme değil bir tahmindir --- sistemini bu gerçeğin üzerine kur, ona rağmen değil.
\end{principlebox}
